\section{Estratégias de escalabilidade}
\subsection{Automação e tecnologia}
Padronizar primeiro e automatizar depois. Priorizar: onboarding de clínicas; geração de QR; lembretes; painel de indicadores; suporte com base de conhecimento; conciliação financeira somente após validação dos fluxos. Integrações externas e infraestrutura remota devem ser testadas em ambiente controlado antes de dados pessoais reais. A escalabilidade técnica será acompanhada por custo de infraestrutura por usuário ativo, disponibilidade, erros, tempo de resposta e incidentes.

\subsection{Modelo replicável}
Criar um ``kit de praça'' composto por checklist de abertura, contrato-padrão revisado juridicamente, critérios de seleção de clínicas, roteiro de parceria com ONGs, materiais QR, treinamento, política de créditos e painel de indicadores. Franquia e licenciamento ficam adiados até existir operação comprovada em mais de uma praça, suporte padronizado e governança de marca.

\subsection{Receita recorrente e salvaguardas}
A receita recorrente poderá vir de PetCuida+, Plano Parceiro e convênios anuais. Não se deve pressupor que todos os parceiros pagarão ao término do período gratuito. O indicador central será receita recorrente efetivamente recebida, com taxa de cancelamento e margem líquida. O fundo de créditos deverá ter teto, regras transparentes, patrocinador identificado e conciliação auditável.